\thispagestyle{empty}

\chapter*{Resumen}

La contaminación atmosférica es uno de los principales retos ambientales y de salud pública, responsable de millones de muertes prematuras y de graves impactos en los ecosistemas y el clima. Las previsiones fiables son fundamentales para las políticas públicas, la concienciación social y la toma de decisiones cotidianas. El servicio Copernicus Atmosphere Monitoring Service (CAMS) ofrece pronósticos de alta calidad a escala global, pero las plataformas actuales presentan una limitación: las herramientas científicas requieren conocimientos técnicos avanzados y las aplicaciones dirigidas al público simplifican en exceso la información.

Este trabajo presenta una plataforma web de visualización que cubre este vacío, ofreciendo una interfaz accesible pero científicamente rigurosa para explorar los pronósticos de CAMS. El sistema funciona íntegramente en el navegador y utiliza tecnologías abiertas como Leaflet.js para mapas interactivos y Highcharts para la visualización de gráficos. Entre sus principales funciones destacan el control deslizante para previsiones horarias, la extracción de valores de contaminantes en cualquier punto, gráficos históricos interactivos, leyendas de color adaptadas a cada contaminantes y la animación de la evolución de los pronósticos.

Los resultados muestran que la plataforma combina con éxito usabilidad y capacidad analítica, permitiendo interpretar previsiones complejas sin necesidad de software especializado. Al estandarizar unidades y añadir leyendas adaptativas, garantiza coherencia y claridad en la interpretación. Además de su valor académico, la herramienta puede apoyar la educación, la concienciación pública y la monitorización ambiental. Como trabajo futuro se plantea integrar datos observacionales para validación e incorporar análisis estadísticos avanzados.
